\section{Structural Reduction to Finite Cases}
\label{sec:structural-reductions}

Throughout, indices are taken in $\mathbb Z/6\mathbb Z$.  Put
\[
 O=(0,0),\qquad
 V_i=\left(\cos\frac{i\pi}{3},\sin\frac{i\pi}{3}\right),
 \qquad
 H=\mathrm{conv}\{V_0,\ldots,V_5\}.
\]
Thus $H$ is the regular hexagon of side length one.  We use the notation
\[
 e_{i,i+1}=[V_i,V_{i+1}],\qquad
 r_i=[O,V_i],\qquad M_i=\frac12V_i,
\]
and, when needed,
\[
 D=\bigcup_{i=0}^5r_i,\qquad
 S=\partial H\cup D,\qquad
 S_{1/2}=\partial H\cup\{O,M_0,\ldots,M_5\}.
\]
The arms $r_i$ meet only at $O$, up to sets of one-dimensional measure zero,
and consequently
\begin{equation}
 \mathcal H^1(\partial H)=6,
 \qquad \mathcal H^1(D)=6,
 \qquad \mathcal H^1(S)=12.
 \label{eq:target-lengths}
\end{equation}
For $L>0$, $H_L:=LH$ denotes the concentric regular hexagon of side length
$L$.  All symmetries used below belong to the dihedral group $D_6$; applying
a symmetry always transports the indices, orientations, and incoming and
outgoing labels together.

\subsection{Open, closed, and scaled formulations}

\begin{proposition}[Compactness and scaling equivalence]
\label{prop:open-closed-scaled}
The following assertions are equivalent.
\begin{enumerate}
 \item The set $H$ is covered by seven open equilateral triangles of side
 length $1$.
 \item For some $\varepsilon>0$, the set $H$ is covered by seven closed
 equilateral triangles of side length $1-\varepsilon$.
 \item For some $L>1$, the set $H_L$ is covered by seven closed equilateral
 triangles of side length $1$.
\end{enumerate}
\end{proposition}

\begin{proof}
Suppose that open unit triangles $U_1,\ldots,U_7$ cover $H$.  Define
\[
 m_j(p)=\mathrm{dist}(p,\mathbb R^2\setminus U_j),
 \qquad m(p)=\max_{1\le j\le7}m_j(p).
\]
The continuous function $m$ is positive on the compact set $H$, so
$\eta:=\min_{p\in H}m(p)>0$.  Choose
\[
 0<\rho<\min\left\{\eta,\frac{\sqrt3}{6}\right\}.
\]
Moving each side of $U_j$ inward through distance $\rho$ gives a closed
equilateral triangle $K_j$ of side $1-2\sqrt3\rho$.  Every $p\in H$ has
$m_j(p)\ge\eta>\rho$ for some $j$, and hence belongs to $K_j$.  This proves
(1)$\Rightarrow$(2), with $\varepsilon=2\sqrt3\rho$.

Conversely, enlarging a closed triangle of side $1-\varepsilon$ about its
incenter to an open triangle of side $1$ preserves containment, proving
(2)$\Rightarrow$(1).  Finally, dilation by $(1-\varepsilon)^{-1}$ identifies
(2) with (3), with
\[
 L=(1-\varepsilon)^{-1}>1.
\]
\end{proof}

We conduct the geometric argument in the open side-one model.  Denote the
original open triangles by $U_C$ and $U_i$ for $0\le i\le5$, and put
\[
 T_C=\overline{U_C},\qquad T_i=\overline{U_i}.
\]
The distinction between $U_i$ and $T_i$ will be maintained: closure may add
endpoints or tangencies, although it does not change area or one-dimensional
trace measure.

\begin{lemma}[Distinct triangles]
\label{lem:distinct-roles}
The seven triangles can be labelled so that
\[
 O\in U_C,\qquad V_i\in U_i\quad(0\le i\le5),
\]
and these seven triangles are distinct.  Consequently
\[
 O\in\mathrm{int}(T_C),\qquad V_i\in\mathrm{int}(T_i).
\]
\end{lemma}

\begin{proof}
The seven distinguished points $O,V_0,\ldots,V_5$ are pairwise at distance
at least $1$.  Any two points of an open unit equilateral triangle are at
distance strictly less than its diameter $1$.  Thus one open triangle
contains at most one distinguished point.  Selecting a covering triangle for
each of the seven points gives an injection from seven points to seven
triangles, hence a bijection and the asserted labelling.
\end{proof}

\subsection{The center classification}

For a closed C triangle define
\[
 N_E(T_C)=
 \left\lvert\left\{i:T_C\cap e_{i,i+1}\text{ contains a
 positive-length interval}\right\}\right\rvert.
\]
Point contacts, endpoint contacts, and tangencies do not contribute.  We say
that $T_C$ has type CE0, CE1, or CE2 according as $N_E(T_C)$ is $0$, $1$, or
$2$.

\begin{proposition}[Exhaustive center types]
\label{prop:ce-classification}
Every C triangle has exactly one of the types CE0, CE1, and CE2.
Equivalently,
\[
 N_E(T_C)\le2.
\]
\end{proposition}

\begin{proof}
Among any three edges of the six-cycle $\partial H$, two have no common
endpoint.  It is therefore enough to show that relative-interior points of
two nonadjacent edges are more than unit distance apart.  Up to symmetry
there are two cases.  For $0<t,s<1$, put
\[
 x=(1-t)V_0+tV_1.
\]
If $y=(1-s)V_2+sV_3$, direct calculation gives
\[
 \lVert x-y\rVert^2-1=(1-t)(2-t)+s(s+t)>0.
\]
If instead $y=(1-s)V_3+sV_4$ and $u=t+s$, then
\[
 \lVert x-y\rVert^2-1=(u-1)^2+2>0.
\]
A unit equilateral triangle has diameter $1$.  It therefore cannot contain
relative-interior points on two nonadjacent hexagon edges.  Three
positive-length edge traces would supply such a pair, a contradiction.
\end{proof}

For the CE1/CE2 types, Proposition~\ref{prop:unique-center-midpoint} below
records the additional fact that $T_C$ contains exactly one radial midpoint.

\subsection{Vertex types and exact-trace normalizations}

For the closure $T_i$ of an original open V triangle $U_i$, let
\[
 o(T_i)=\left\lvert\left\{
 \text{vertices of }T_i\text{ outside }H
 \right\}\right\rvert
\]
and put
\[
 n(T_i)=\left\lvert\left\{j\in\{i-1,i+1\}:
 \Haus(T_i\cap r_j)>0\right\}\right\rvert.
\]

\begin{lemma}[Local wedge]
\label{lem:local-wedge}
Every original V triangle has $1\le o\le3$.  If
$\Haus(T_i\cap r_j)>0$ for some $j\in\{i-1,i+1\}$, at least one of its
vertices belongs to $H$; if both $\Haus(T_i\cap r_{i-1})>0$ and
$\Haus(T_i\cap r_{i+1})>0$, at least two of its vertices belong to $H$.
\end{lemma}

\begin{proof}
Normalize to $i=0$ and write
\[
 X=V_0+x(V_1-V_0)+y(V_5-V_0).
\]
Then $V_0=(0,0)$,
\[
 \lVert(x,y)\rVert^2=x^2+y^2-xy,
\]
and, within the closed unit ball about $V_0$, membership in $H$ is equivalent
to membership in the wedge
\[
 W=\{(x,y):x\ge0, y\ge0\}.
\]
Indeed,
\[
 x^2+y^2-xy=\left(y-\frac x2\right)^2+\frac{3x^2}{4}
              =(x-y)^2+xy
\]
shows that $x,y\ge0$ and norm at most one imply the remaining hexagon
inequalities.  The two relevant radial segments are
\[
 r_1=\{(1,y):0\le y\le1\},
 \qquad r_5=\{(x,1):0\le x\le1\}.
\]

Suppose $T_0\cap r_1$ has positive length, and let $m$ be the largest
$x$-coordinate of a vertex of $T_0$.  If $m>1$, a maximizing vertex
$(m,y)$ lies in the unit ball and hence satisfies $0<y<1$; it belongs to
$H$.  If $m=1$, positive-length contact with the supporting line $x=1$
forces a side of $T_0$ to lie there, and both endpoints satisfy
$1+y^2-y\le1$, hence belong to $H$.  The same argument applies to $r_5$.
If both $r_1$ and $r_5$ are met but only one triangle vertex belonged to $H$, that
vertex would have to have both $x>1$ and $y>1$.  This is impossible because
$(x-y)^2+xy>1$.

Finally, if all three triangle vertices belonged to the convex set $H$, then
$T_0\subset H$.  This is incompatible with the extreme point $V_0$ lying in
the interior of $T_0$.  Hence $o\ge1$.
\end{proof}

We use the following orientation normal forms for both translations below.
Set
\[
 D_t=1-t+t^2,\qquad z=\sqrt{D_t},\qquad 0\le t\le1.
\]
After relabelling the vertices of the equilateral triangle, every orientation
belongs to one of the two families
\[
 \begin{array}{c|cc}
 &\text{first side vector}&\text{second side vector}\\ \hline
 \mathrm{I}&z^{-1}(1,t)&z^{-1}(1-t,1)\\
 \mathrm{II}&z^{-1}(t,-(1-t))&z^{-1}(1,t).
 \end{array}
\]
In Type I write
\begin{equation}
 U=\alpha+y-tx,\qquad V=\beta+x-(1-t)y,
 \qquad T=\{U\ge0,V\ge0,U+V\le z\},
 \label{eq:orientation-type-I}
\end{equation}
and in Type II write
\begin{equation}
 U=\alpha+(1-t)x+ty,\qquad V=\beta+tx-y,
 \qquad T=\{U\ge0,V\ge0,U+V\le z\}.
 \label{eq:orientation-type-II}
\end{equation}
In either form, $V_0\in\interior(T)$ is equivalent to
\begin{equation}
 \alpha>0,\qquad\beta>0,\qquad\alpha+\beta<z.
 \label{eq:orientation-interior}
\end{equation}
For later use, the Type-II vertex in the wedge is
\begin{equation}
 Q=\left(
 \frac{z-\alpha-t\beta}{D_t},
 \frac{t(z-\alpha)+(1-t)\beta}{D_t}
 \right).
 \label{eq:orientation-type-II-Q}
\end{equation}
Both coordinates are positive under \eqref{eq:orientation-interior}.

\begin{proposition}[Exact-trace normalization of raw $(3,0)$ roles]
\label{prop:vd0-exact-trace-normalization}
Let $U$ be an original open $V_i$-triangle, put $T=\overline U$, and suppose
$(o(T),n(T))=(3,0)$.  There is a translate $\widehat U$ of $U$, with closure
$\widehat T$, the same orientation, and the same side length, such that
\begin{equation}
 V_i\in\widehat U,
 \qquad
 (o(\widehat T),n(\widehat T))=(2,0),
 \label{eq:o3-normalized-type}
\end{equation}
and both the closed and open traces on the hexagon are unchanged:
\begin{equation}
 T\cap H=\widehat T\cap H,
 \qquad
 U\cap H=\widehat U\cap H.
 \label{eq:o3-exact-traces}
\end{equation}
Consequently the three maximal reaches on the two incident edges and the own
radial segment are unchanged.  In the notation introduced below, the values
$A_i,B_i,C_i$, the inequality $A_i+B_i>1$, and hence $N_+$ are all preserved.
\end{proposition}

\begin{proof}
Normalize to $i=0$ and use the common forms above.  A Type-II triangle cannot
have $o=3$: the point $Q$ in \eqref{eq:orientation-type-II-Q} lies in $W$,
and because $T$ contains the origin and has diameter one, the local wedge
lemma puts $Q$ in $H$.  Thus $T$ has Type I form.

The three Type-I vertices are
\[
 Q_0=\left(
 -\frac{(1-t)\alpha+\beta}{D_t},
 -\frac{\alpha+t\beta}{D_t}
 \right),
\]
\begin{equation}
 \begin{aligned}
 Q_1&=\left(
 \frac{z-(1-t)\alpha-\beta}{D_t},
 \frac{tz-\alpha-t\beta}{D_t}
 \right),\\
 Q_2&=\left(
 \frac{(1-t)z-(1-t)\alpha-\beta}{D_t},
 \frac{z-\alpha-t\beta}{D_t}
 \right).
 \end{aligned}
 \label{eq:o3-type-I-vertices}
\end{equation}
The vertex $Q_0$ has two negative coordinates, the first coordinate of $Q_1$
is positive, and the second coordinate of $Q_2$ is positive.  Hence $o=3$
is equivalent to
\begin{equation}
 \alpha+t\beta>tz,
 \qquad
 (1-t)\alpha+\beta>(1-t)z.
 \label{eq:o3-outside-conditions}
\end{equation}
These strict inequalities force $0<t<1$.

Put
\[
 s=\alpha+\beta,
 \qquad L=z-s>0.
\]
The two inequalities in \eqref{eq:o3-outside-conditions} become
\begin{equation}
 \alpha>\frac{tL}{1-t},
 \qquad
 \beta>\frac{(1-t)L}{t}.
 \label{eq:o3-redundancy-conditions}
\end{equation}
For every $(x,y)\in W$ satisfying
\begin{equation}
 (1-t)x+ty\le L,
 \label{eq:o3-wedge-cut}
\end{equation}
the minimum of $U$ is attained at $(L/(1-t),0)$ and the minimum of $V$ at
$(0,L/t)$.  Thus \eqref{eq:o3-redundancy-conditions} makes both $U$ and $V$
strictly positive throughout the triangle in \eqref{eq:o3-wedge-cut}.  It
follows that
\begin{equation}
 T\cap W=\left\{(x,y)\in W:(1-t)x+ty\le L\right\}.
 \label{eq:o3-wedge-trace}
\end{equation}

Keep $t$ and $s$ fixed and replace the affine constants by
\begin{equation}
 \widehat\alpha=\frac{tL}{1-t},
 \qquad
 \widehat\beta=s-\widehat\alpha.
 \label{eq:o3-normalized-constants}
\end{equation}
The first inequality in \eqref{eq:o3-redundancy-conditions} gives
$0<\widehat\alpha<\alpha$ and
$\widehat\beta>\beta>0$; also
$\widehat\alpha+\widehat\beta=s<z$.  Since the linear map defining $U,V$ in
\eqref{eq:orientation-type-I} is nonsingular, changing only these affine
constants gives a translate $\widehat T$ of $T$.  In particular,
\eqref{eq:orientation-interior} shows that the
origin remains in $\interior(\widehat T)$.

For the translated triangle, the inequality $\widehat U+\widehat V\le z$ is
still exactly \eqref{eq:o3-wedge-cut}.  On that triangle $\widehat U\ge0$,
with equality only at $(L/(1-t),0)$, while $\widehat V>0$ by the second
inequality in \eqref{eq:o3-redundancy-conditions} and
$\widehat\beta>\beta$.  Therefore
\[
 \widehat T\cap W=T\cap W,
\]
and their interiors have the same trace on $W$: the only new equality point
of $\widehat U$ already lies on the common boundary
$(1-t)x+ty=L$.  Both triangles contain the origin and have diameter one, so
the local wedge lemma identifies their intersections with $H$ with their
intersections with $W$.  This proves both equalities in
\eqref{eq:o3-exact-traces}.

Equation \eqref{eq:o3-normalized-constants} makes the second coordinate of
$Q_1$ in \eqref{eq:o3-type-I-vertices} zero.  Thus $Q_1\in H$.  The vertex
$Q_0$ remains outside $H$, and the strict second inequality in
\eqref{eq:o3-redundancy-conditions}, strengthened by
$\widehat\beta>\beta$, keeps $Q_2$ outside $H$.  Hence
$o(\widehat T)=2$.  The local wedge lemma gives $n(T)=0$ from $o(T)=3$, and
the exact trace equality preserves $n$, so $n(\widehat T)=0$.  Finally, all
three reach segments lie in $H$, and \eqref{eq:o3-exact-traces} preserves
their exact open and closed traces.  This proves every assertion.
\end{proof}

Apply Proposition~\ref{prop:vd0-exact-trace-normalization} simultaneously to
every raw $(3,0)$ vertex role and then reuse the symbols $U_i,T_i$ for the
translated roles.  The rolewise equalities \eqref{eq:o3-exact-traces}
preserve the entire open cover, every closed trace, the condition
$V_i\in U_i$, all actual reaches, and $N_+$.  This simultaneous exact-trace
normalization is the convention throughout the rest of the paper.

For the normalized roles, we use the names
\[
 \begin{array}{c|c}
 \text{type}&(o,n)\\ \hline
 \mathrm{Vd0}&(1,0),(2,0)\\
 \mathrm{Vd1}&(1,1)\\
 \mathrm{Vd2}&(1,2)\\
 \text{T3-like}&(2,1).
 \end{array}
\]

\begin{proposition}[Exhaustive normalized vertex types]
\label{prop:vertex-classification}
Every normalized V triangle has exactly one of the four types Vd0, Vd1, Vd2,
and T3-like.
\end{proposition}

\begin{proof}
We have $n\in\{0,1,2\}$.  Lemma~\ref{lem:local-wedge} gives $o\le2$ when
$n\ge1$, and $o\le1$ when $n=2$.  Together with $o\ge1$, the possibilities
for a raw role are
\[
 \begin{array}{c|c}
 n&\text{possible }o\\ \hline
 0&1,2,3\\
 1&1,2\\
 2&1.
 \end{array}
\]
The exact-trace normalization replaces the sole extra possibility $(3,0)$ by
$(2,0)$ and leaves all other pairs unchanged.  The remaining pairs are
precisely the four named classes.
\end{proof}

The next normalization is useful for trace estimates, but its final warning
is essential.

\begin{proposition}[T3-like closed-trace domination]
\label{prop:t3-translation}
Let $T$ be the closure of an original T3-like $V_i$-triangle.  There is a
translate $\widehat T$ of $T$, with the same orientation and side length,
such that
\[
 V_i\in\mathrm{relint}(\text{a side of }\widehat T),
 \qquad T\cap H\subsetneq\widehat T\cap H,
\]
and $\widehat T$ remains of type $(o,n)=(2,1)$.  Every old open interior
point in $H\setminus\{V_i\}$ remains an interior point of $\widehat T$.
The point $V_i$ itself becomes a boundary point, so $\widehat T$ is a
closed-trace majorant and is not a replacement open triangle.
\end{proposition}

\begin{proof}
Use the wedge coordinates of Lemma~\ref{lem:local-wedge} and the common
orientation normal forms above.  In Type I, direct inversion shows that if
exactly two triangle vertices are outside $W$, the sole wedge vertex has
both coordinates strictly below $1$: in the two possible cases its
coordinates are bounded respectively by $(1-t)/z$ and $t/z$, with the
endpoint strictness supplied by $\alpha>0$ or $\beta>0$.  The support
argument in Lemma~\ref{lem:local-wedge} then excludes positive-length contact
with either $x=1$ or $y=1$.  Thus Type I cannot be T3-like.

The triangle therefore has Type II form.
Reflecting in $x=y$ if necessary, assume $\Haus(T_0\cap r_1)>0$.  The two
outside vertices have respectively a negative $x$- and a
negative $y$-coordinate, while the unique wedge vertex is the point $Q$ in
\eqref{eq:orientation-type-II-Q}.
The endpoint cases $t=0,1$ give no positive adjacent interval, so
$0<t<1$.  The exact axis reaches are
\[
 A=\beta,\qquad B=z-\alpha-\beta,
\]
and T3-like support is equivalent to
\[
 Q_x>1,\qquad Q_y\le1.
\]

Define $\widehat T$ by replacing $\alpha$ with $0$ and keeping $t,\beta$
fixed.  Since the linear map $(x,y)\mapsto(U,V)$ is nonsingular, this changes
only the translation.  At the origin, $\widehat U=0$ and
$0<\widehat V=\beta<z$, so $V_0$ is in the relative interior of a side.  For
$(x,y)\in W$,
\[
 \widehat U=(1-t)x+ty\ge0,\qquad
 \widehat V=V,\qquad
 \widehat U+\widehat V=U+V-\alpha.
\]
Thus $T\cap W\subset\widehat T\cap W$.  The old $x$-axis reach is $B$ and
the new one is $B+\alpha$; because $B<1$ and $\alpha>0$, the inclusion is
strict inside $H$.

The two outside vertices remain outside, whereas $V_0$ and the new wedge
vertex lie in $H$, so $o=2$.  Moreover $\widehat Q_x>Q_x>1$.  From $Q_x>1$
one gets $t\beta<z-\alpha-D_t$, and hence
\[
 tD_t(1-\widehat Q_y)
 >D_t(1-z)+(1-t)\alpha>0.
\]
Thus $\widehat Q_y<1$ and $n=1$.  Finally, for every nonzero point of $W$,
$\widehat U>0$; the displayed slack relations show that every old interior
point other than the origin remains interior.
\end{proof}

\subsection{Reach coordinates, supercritical V triangles, and gaps}

For the V triangle $U_i$ under the simultaneous normalization convention,
define its maximal incoming, outgoing, and radial reaches by
\[
 \begin{aligned}
 A(U_i)&=
 \sup_{\substack{t\in[0,1]\\V_i+t(V_{i-1}-V_i)\in U_i}}t,\\
 B(U_i)&=
 \sup_{\substack{t\in[0,1]\\V_i+t(V_{i+1}-V_i)\in U_i}}t,\\
 C(U_i)&=
 \sup_{\substack{t\in[0,1]\\V_i+t(O-V_i)\in U_i}}t.
 \end{aligned}
\]
Closure does not change the suprema, so
\[
 A(U_i)=A(T_i),\qquad B(U_i)=B(T_i),\qquad C(U_i)=C(T_i).
\]
For compact indexed formulas, put
\[
 A_i:=A(U_i)=A(T_i),\qquad
 B_i:=B(U_i)=B(T_i),\qquad
 C_i:=C(U_i)=C(T_i).
\]
The closure $T_i$ contains the corresponding endpoints. We reserve
lowercase $(a_i,b_i,c_i)$ for selected or prescribed reach lower bounds
whenever actual and selected values occur together.  A realizing triangle then
satisfies
\[
 A_i\ge a_i,\qquad B_i\ge b_i,\qquad C_i\ge c_i.
\]
An actual V triangle is supercritical if $A_i+B_i>1$, and
\begin{equation}
 N_+=\left\lvert\left\{i:A_i+B_i>1\right\}\right\rvert.
 \label{eq:N-plus-definition}
\end{equation}
This definition never counts a smaller reach pair in place of the actual
V triangle.

Parametrize $e_{i,i+1}$ by
\[
 X_i(x)=V_i+x(V_{i+1}-V_i),\qquad 0\le x\le1.
\]
The two incident open traces are
\[
 U_i\cap e_{i,i+1}=[0,B_i),
 \qquad
 U_{i+1}\cap e_{i,i+1}=(1-A_{i+1},1]
\]
in this coordinate.  A \emph{boundary gap} on this edge is the nonempty
closed complement between these two traces,
\[
 [B_i,1-A_{i+1}]\qquad\text{when }B_i\le1-A_{i+1}.
\]
In particular, equality gives a singleton gap.  This convention retains the
actual open traces; a point missing from both open triangles is not erased merely
because both closures contain it.

\begin{lemma}[Exhaustiveness of the gap split]
\label{lem:gap-exhaustion}
On each edge the complement of the two incident V triangle traces is either
empty or one closed interval; a singleton interval counts as a gap.  Every
edge having a gap has positive-length intersection with the open center
triangle.  Consequently a $\CEzero$ center permits no gaps, a $\CEone$ center
permits at most one, and a $\CEtwo$ center permits at most two.
\end{lemma}

\begin{proof}
The two incident traces are intervals issuing from the two endpoints, so
their complement is exactly the interval displayed above when nonempty.  A
nonincident V triangle cannot contain a relative-interior point of the edge:
that point is more than unit distance from the triangle's distinguished vertex.
Thus every point of a gap lies in $U_C$.  Because $U_C$ is open, its
intersection with that edge contains a relative neighborhood and hence has
positive length.  Different gap edges therefore give different edges counted
by $N_E(T_C)$.  Proposition~\ref{prop:ce-classification} gives the three
stated bounds, including singleton gaps.
\end{proof}

\begin{lemma}[T3-like V triangles are nonsupercritical]
\label{lem:t3-nonsupercritical}
Every original T3-like V triangle satisfies
\begin{equation}
 A_i+B_i\le1.
\label{eq:t3-nonsupercritical}
\end{equation}
In particular, it is not counted by $N_+$.
\end{lemma}

\begin{proof}
The proof of Proposition~\ref{prop:t3-translation} showed that a T3-like
triangle has Type II orientation with
\[
 A_i=\beta,\qquad B_i=z-\alpha-\beta,\qquad
 z=\sqrt{1-t+t^2}\le1.
\]
Thus $A_i+B_i=z-\alpha\le1$.  The same conclusion holds after reflection.
\end{proof}

We state the precise passage from maximal reaches to strict handoff lower bounds.
Its hypothesis that the six V triangles themselves cover the boundary is
part of the statement.

\begin{proposition}[Strict boundary handoffs]
\label{prop:strict-handoffs}
Assume that $U_0,\ldots,U_5$ cover $\partial H$.  On each edge choose
\[
 \xi_i\in(1-A_{i+1},B_i)
\]
and define
\[
 a_i=1-\xi_{i-1},\qquad b_i=\xi_i.
\]
Then $0<a_i<A_i$, $0<b_i<B_i$, the two selected anchors lie in $U_i$, and
\[
 a_i^2+a_ib_i+b_i^2<1.
\]
Moreover:
\begin{enumerate}
 \item if exactly one actual V triangle, indexed by $p$, is supercritical, every
 strict selection has exactly one selected supercritical V triangle, also indexed
 by $p$;
 \item if at least two actual V triangles are supercritical, some simultaneous
 strict selection has at least two selected supercritical V triangles.
\end{enumerate}
\end{proposition}

\begin{proof}
Because $V_i$ is interior to a diameter-one closed triangle,
\[
 0<A_i<1,\qquad 0<B_i<1.
\]
No nonincident V triangle contains a relative-interior point of
$e_{i,i+1}$, since such a point is more than unit distance from its
distinguished vertex.  The two displayed incident traces therefore cover
the edge and, being relatively open and nonempty, overlap.  Hence each
interval $(1-A_{i+1},B_i)$ is nonempty.  The anchor and distance assertions
follow from the choice and from the fact that two points of $U_i$ are at
distance strictly less than one.

Put $\ell_i=1-A_{i+1}$ and $u_i=B_i$.  Then
\begin{equation}
 A_i+B_i>1\iff \ell_{i-1}<u_i,
 \qquad
 a_i+b_i>1\iff \xi_{i-1}<\xi_i.
 \label{eq:actual-selected-ascents}
\end{equation}
If actual V triangle $i$ is nonsupercritical, then
\[
 \xi_i<u_i\le\ell_{i-1}<\xi_{i-1},
\]
so every selected version of it is strictly subcritical.  Also
\begin{equation}
 \sum_{i=0}^5(a_i+b_i)
 =\sum_{i=0}^5(1-\xi_{i-1}+\xi_i)=6.
 \label{eq:handoff-telescoping}
\end{equation}
If $p$ is the sole actual supercritical index, the other five selected sums
are strictly below one; equation~\eqref{eq:handoff-telescoping} forces the
$p$th sum above one.  This proves (1).

For (2), first take two nonadjacent supercritical indices $p,q$.  For each
$r\in\{p,q\}$ choose
\[
 \ell_{r-1}<z_r<u_r,
\]
then choose
\[
 \xi_{r-1}\in(\ell_{r-1},\min\{u_{r-1},z_r\}),
 \quad
 \xi_r\in(\max\{\ell_r,z_r\},u_r).
\]
The two pairs of variables are disjoint and give ascents at $p,q$.  Any set
of at least three indices on a six-cycle contains two nonadjacent indices.
It remains to consider exactly two adjacent supercritical indices $p,p+1$.
For the other four nonsupercritical V triangles, successive inequalities
$\ell_i<u_i\le\ell_{i-1}$ give
\[
 \ell_{p-1}<\ell_{p+1}.
\]
Together with supercriticality at $p,p+1$, this implies
\[
 \max\{\ell_{p-1},\ell_p\}<\min\{u_p,u_{p+1}\}.
\]
Choose $\xi_p$ between these bounds and then choose
$\xi_{p-1}<\xi_p<\xi_{p+1}$ in their respective overlap intervals.  By
\eqref{eq:actual-selected-ascents}, both $T_p$ and $T_{p+1}$ are selected
supercritical.
\end{proof}

\subsection{Midpoint forcing}

We first need a local fact for V triangles.  Take unit vectors
\[
 u=\left(\frac12,\frac{\sqrt3}{2}\right),\qquad
 v=\left(\frac12,-\frac{\sqrt3}{2}\right).
\]

\begin{lemma}[Self-midpoint obstruction]
\label{lem:self-midpoint}
Let a closed unit triangle contain
\[
 0,\qquad au,\qquad bv,\qquad \frac{u+v}{2}.
\]
Then $a+b\le1$.  Consequently, for every V triangle,
\[
 M_i\in T_i\quad\Longrightarrow\quad A_i+B_i\le1.
\]
If the two demanded boundary points and the midpoint are contained with a
positive margin, the conclusion is strict.
\end{lemma}

\begin{proof}
Suppose $a+b>1$ and set
\[
 K=\mathrm{conv}\left\{0,au,\frac{u+v}{2},bv\right\}.
\]
For outward unit normals $n_0,n_1,n_2$ separated by $120^\circ$, the least
equilateral triangle with those normals and containing $K$ has side length
\[
 \frac2{\sqrt3}\sum_{j=0}^2h_K(n_j).
\]
The finite-caliper theorem, Theorem~\ref{thm:cert-caliper}, reduces its
minimum to an outward normal of one of the four edges of $K$.

For the edge from $0$ to $au$, direct projection gives
\[
 S_{0A}=\max\left\{\frac{\sqrt3}{4},\frac{\sqrt3 a}{2}\right\}
          +\frac{\sqrt3 b}{2}>\frac{\sqrt3}{2};
\]
if $a\ge1/2$ this is $(\sqrt3/2)(a+b)$, and if $a<1/2$ then
$b>1/2$.  The edge through $bv$ is symmetric.  For the edge from $au$ to
$(u+v)/2$, put $D_a=\sqrt{4a^2-2a+1}$.  Projection gives
\[
 S_{AM}=
 \begin{cases}
 \dfrac{\sqrt3(a+b)}{2D_a},&a\le\frac12,\\[4pt]
 \dfrac{\sqrt3(2a^2+b)}{2D_a},&a>\frac12.
 \end{cases}
\]
The first value is strictly above $\sqrt3/2$ because $D_a\le1$.  In the
second case, $b>1-a$ and
\[
 (2a^2-a+1)^2-D_a^2=a^2(2a-1)^2\ge0,
\]
so again $S_{AM}>\sqrt3/2$.  The fourth edge is symmetric.  Every enclosing
equilateral triangle therefore has side greater than one, a contradiction.

If equality $a+b=1$ occurred while all relevant points had positive margin,
both boundary reach lower bounds could be increased slightly; the closed conclusion
would then be violated.
\end{proof}

The corresponding center fact is exact and uses the interior condition at
$O$.

\begin{proposition}[CE1/CE2 exactly-one-midpoint property]
\label{prop:unique-center-midpoint}
Let $T_C$ be a closed unit equilateral triangle with
$O\in\mathrm{int}(T_C)$.  If $T_C$ is CE1 or CE2, equivalently if it has
positive-length intersection with a boundary edge of $H$, then
\[
 \left\lvert T_C\cap\{M_0,\ldots,M_5\}\right\rvert=1.
\]
\end{proposition}

\begin{proof}
Normalize the positive trace by a dihedral symmetry.  Proposition
\ref{prop:signed-center-normal-form} gives a single signed side-slack model
for both CE1 and CE2 and proves directly that its midpoint set is
$\{M_0\}$.  Undoing the normalization proves the assertion.
\end{proof}

\begin{remark}
The hypothesis $O\in\mathrm{int}(T_C)$ cannot be weakened to
$O\in T_C$: the triangle $\mathrm{conv}\{O,V_0,V_1\}$ has a boundary-edge
overlap and contains both $M_0$ and $M_1$.
\end{remark}

\begin{proposition}[Exhaustive structural reduction]
\label{prop:exhaustive-structural-reduction}
Every hypothetical cover by seven open unit equilateral triangles admits the
distinct center and V triangles used in Table~\ref{tab:routing}, and its
maximal reaches determine exactly one row of that table.
\end{proposition}

\begin{proof}
Lemma~\ref{lem:distinct-roles} gives the seven triangles.
Propositions~\ref{prop:ce-classification} and
\ref{prop:vertex-classification} give the exhaustive and mutually exclusive
center and vertex types.  The maximal reaches determine exactly one of
$N_+=0$, $N_+=1$, and $N_+\ge2$.  For CE0 these data and the vertex
classification give the five displayed rows.

Now suppose the center is CE1 or CE2.  It contains exactly one radial
midpoint by Proposition~\ref{prop:unique-center-midpoint}.  The direct
skeleton theorem, Theorem~\ref{thm:common-skeleton-count}, removes every
state with $N_++N_{\mathrm{sp}}\ge3$.  If $N_+\ge2$, Lemma
\ref{lem:positive-support-rescuer} supplies a distinct positive-support V
triangle, so $N_{\mathrm{sp}}\ge1$ and that theorem applies.  Hence every
surviving branch has $N_+\le1$ and $N_++N_{\mathrm{sp}}\le2$.  If $N_+=0$,
the vertex classification gives, in order, all $\Vdzero$, at least one
$\Vdone/\Vdtwo$, or no such triangle and at least one $\Tlike$.  If $N_+=1$,
then $N_{\mathrm{sp}}\le1$: the triangles are all $\Vdzero$, exactly one
$\Tlike$, or exactly one $\Vdone/\Vdtwo$.  In the all-$\Vdzero$ case,
Lemma~\ref{lem:gap-exhaustion} supplies the exhaustive zero-gap/nonempty-gap
split.  These alternatives are disjoint, proving the claim.
\end{proof}
